\section{Magnetic-force-theorem mixed derivative}
\label{sec:mixed-response}

At fixed chemical potential, the one-particle free-energy derivative can be written as
\begin{equation}
\frac{\partial\bF}{\partial\lambda}
=-\frac{1}{\pi}\operatorname{Im}
\int\dd\varepsilon\,f(\varepsilon)
\Tr[V_\lambda G^R],
\qquad
V_\lambda=\frac{\partial H}{\partial\lambda}.
\label{eq:first-derivative}
\end{equation}
Using
\begin{equation}
\frac{\partial G}{\partial\eta}=GV_\eta G,
\label{eq:dyson-derivative}
\end{equation}
we obtain
\begin{equation}
\frac{\partial^2\bF}{\partial\lambda\partial\eta}
=-\frac{1}{\pi}\operatorname{Im}
\int\dd\varepsilon\,f(\varepsilon)
\Tr\left[V_{\lambda\eta}G^R+V_\lambda G^R V_\eta G^R\right].
\label{eq:mixed-general}
\end{equation}
The second term is the mixed bubble. The same perturbative structure appears in the site-diagonal spin-lattice torque derivation of Ref.~\cite{Mankovsky2023SLC}; there the propagator and lattice vertex are KKR objects.

For a DFPT displacement,
\begin{equation}
g_{\kappa\mu}(\kk,\qq)
=\left\langle w_{\kk+\qq}\middle|
\frac{\partial H}{\partial u_{\kappa\mu}(\qq)}
\middle|w_{\kk}\right\rangle.
\label{eq:dfpt-g}
\end{equation}
This is the screened Hamiltonian derivative supplied by DFPT and transformed into the same Wannier gauge as $H$. The bubble-only approximation sets the mixed Hamiltonian derivative $V_{\lambda\eta}$ to zero. This is the minimal independent-perturbation implementation; the direct term is discussed in \cref{app:direct}.

The target energy coupling is
\begin{equation}
\mathcal H_{\rm SL}^{(1,1)}
=\sum_{\qq,\ell a,\kappa\mu}
\pi_{\ell a}(-\qq)
K_{\ell a,\kappa\mu}(\qq)
 u_{\kappa\mu}(\qq).
\label{eq:coupling-energy}
\end{equation}
The corresponding tangent effective field is $H^{\rm eff}_{\ell a}=-K_{\ell a,\kappa\mu}u_{\kappa\mu}$, while the physical torque is $\bm\tau_\ell=\bm n_\ell\times\bm H^{\rm eff}_\ell$, consistent with the torque-energy relation summarized in Ref.~\cite{Mankovsky2023SLC}.
